% Written by Anders Sjoqvist and Ulf Lundstrom, 2009
% The main sources are: tinyKACTL, Beta and Wikipedia

\chapter{Mathematics}

\section{Sums}
\[ c^a + c^{a+1} + \dots + c^{b} = \frac{c^{b+1} - c^a}{c-1}, c \neq 1 \]
\begin{align*}
	1 + 2 + 3 + \dots + n &= \frac{n(n+1)}{2} \\
	1^2 + 2^2 + 3^2 + \dots + n^2 &= \frac{n(2n+1)(n+1)}{6} \\
	1^3 + 2^3 + 3^3 + \dots + n^3 &= \frac{n^2(n+1)^2}{4} \\
	1^4 + 2^4 + 3^4 + \dots + n^4 &= \frac{n(n+1)(2n+1)(3n^2 + 3n - 1)}{30} \\
\end{align*}

\section{Probability theory}
Let $X$ be a discrete random variable with probability $p_X(x)$ of assuming the value $x$. It will then have an expected value (mean) $\mu=\mathbb{E}(X)=\sum_xxp_X(x)$ and variance $\sigma^2=V(X)=\mathbb{E}(X^2)-(\mathbb{E}(X))^2=\sum_x(x-\mathbb{E}(X))^2p_X(x)$ where $\sigma$ is the standard deviation. If $X$ is instead continuous it will have a probability density function $f_X(x)$ and the sums above will instead be integrals with $p_X(x)$ replaced by $f_X(x)$.

Expectation is linear:
\[\mathbb{E}(aX+bY) = a\mathbb{E}(X)+b\mathbb{E}(Y)\]
For independent $X$ and $Y$, \[V(aX+bY) = a^2V(X)+b^2V(Y).\]

\subsection{Discrete distributions}

\subsubsection{Binomial distribution}
The number of successes in $n$ independent yes/no experiments, each which yields success with probability $p$ is $\textrm{Bin}(n,p),\,n=1,2,\dots,\, 0\leq p\leq1$.
\[p(k)=\binom{n}{k}p^k(1-p)^{n-k}\]
\[\mu = np,\,\sigma^2=np(1-p)\]
$\textrm{Bin}(n,p)$ is approximately $\textrm{Po}(np)$ for small $p$.

\subsubsection{First success distribution}
The number of trials needed to get the first success in independent yes/no experiments, each which yields success with probability $p$ is $\textrm{Fs}(p),\,0\leq p\leq1$.
\[p(k)=p(1-p)^{k-1},\,k=1,2,\dots\]
\[\mu = \frac1p,\,\sigma^2=\frac{1-p}{p^2}\]

\section{Floor sum}
\kactlimport{floorSum.h}
For non-negative integers $a,b,n$ and $c>0$, define
\[F(a,b,c,n)=\sum_{i=0}^{n}\left\lfloor\frac{ai+b}{c}\right\rfloor.\]
Let $T_n=\frac{n(n+1)}{2}$. If $a\ge c$ or $b\ge c$,
\[F(a,b,c,n)=T_n\left\lfloor\frac ac\right\rfloor
+(n+1)\left\lfloor\frac bc\right\rfloor
+F(a\bmod c,b\bmod c,c,n).\]
Otherwise, let $m=\left\lfloor\frac{an+b}{c}\right\rfloor$. Then
\[F(a,b,c,n)=
\begin{cases}
0 & m=0,\\
mn-F(c,c-b-1,a,m-1) & m>0.
\end{cases}\]
The recursion is based on the Euclidean algorithm and runs in
$\mathcal{O}(\log\max(a,c))$ time.

\subsection{Sum of $\lfloor n/i \rfloor$}
To compute
\[S(n)=\sum_{i=1}^{n}\left\lfloor\frac{n}{i}\right\rfloor,\]
group consecutive indices with equal quotients. For each $l$, let
\[q=\left\lfloor\frac{n}{l}\right\rfloor,\qquad
r=\left\lfloor\frac{n}{q}\right\rfloor.\]
Then $\lfloor n/i\rfloor=q$ for all $l\le i\le r$, so add
$q(r-l+1)$ and continue from $l=r+1$.
The algorithm runs in $\mathcal{O}(\sqrt n)$ time.

\section{Linear recurrences}

For a recurrence
\[f_n=c_1f_{n-1}+\dots+c_kf_{n-k},\]
define the state vector
\[v_n=(f_n,f_{n-1},\dots,f_{n-k+1})^T.\]
Then $v_n=Mv_{n-1}$, where
\[
M=\begin{pmatrix}
c_1&c_2&\dots&c_{k-1}&c_k\\
1&0&\dots&0&0\\
0&1&\dots&0&0\\
\vdots&\vdots&\ddots&\vdots&\vdots\\
0&0&\dots&1&0
\end{pmatrix}.
\]
The first row computes the next term; the remaining rows
shift the previous terms. Thus,
\[v_n=M^{n-k+1}v_{k-1},\]
where $v_{k-1}=(f_{k-1},\dots,f_0)^T$.